\subsection{Факторы и институты формирования политической культуры: процесс политической социализации}

Формирование и воспроизводство политической культуры общества представляют собой непрерывный динамический процесс, обозначаемый в политической науке термином «политическая социализация». Под политической социализацией понимается процесс усвоения индивидом ценностей, норм, традиций и образцов поведения, присущих данной политической системе, позволяющих ему адаптироваться к существующему властному порядку и полноценно выполнять гражданские роли. Этот процесс обеспечивает латентную преемственность политической жизни при смене поколений.

Политическая социализация носит многофакторный характер и осуществляется через систему специализированных и неспециализированных социальных институтов (агентов социализации). В отечественной политологической традиции, представленной трудами А.И. Соловьева, агенты социализации классифицируются по степени их формализации и этапам воздействия на личность.

Первичным и наиболее устойчивым институтом формирования политической культуры выступает семья. Именно в семейной микросреде в раннем детстве закладываются базовые, часто неосознанные, эмоционально-оценочные ориентации личности. Ребенок усваивает первые паттерны отношения к авторитету, лидерству, справедливости и порядку. Психологический климат в семье (авторитарный или демократический) предопределяет будущую склонность индивида к подданническому подчинению или активистскому участию. Как отмечают исследователи, ценностные установки, сформированные в семье, обладают наивысшей степенью устойчивости и с трудом поддаются последующей трансформации\footnote{Соловьев А. И. Политология: Политическая теория, политические технологии: Учебник для студентов вузов. --- М.: Аспект Пресс, 2003. --- С. 351.}.

Вторым важнейшим фактором является система образования (школа, средние специальные и высшие учебные заведения). В отличие от семьи, образование выступает как формализованный государственный институт, осуществляющий целенаправленное и рациональное структурирование политического сознания. Здесь происходит трансляция официальной исторической памяти, основ правовых знаний, гражданских обязанностей и государственной символики. Система образования формирует когнитивный (познавательный) компонент политической культуры, переводя первичные семейные эмоции в плоскость осознанных политических убеждений и идеологических предпочтений.

В современном информационном обществе доминирующим фактором конструирования политической культуры становятся средства массовой информации (СМИ) и цифровые коммуникации. Обладая способностью к масштабному и оперативному воздействию, современные медиа фактически формируют виртуальную политическую реальность. Они не просто транслируют информацию, но и задают повестку дня, внедряют определенные интерпретационные рамки (фреймы) и стереотипы восприятия власти и оппозиции. Особое значение приобретает феномен сетевой социализации в интернет-пространстве, который нередко ведет к фрагментации национальной политической культуры на изолированные виртуальные субкультуры.

Дополнительными латентными факторами выступают религиозные институты, фиксирующие традиционные духовно-нравственные основания политики, а также группы сверстников (микроокружение), обеспечивающие неформальное освоение поведенческих моделей. Таким образом, политическая культура нации формируется в результате сложного результирующего взаимодействия всех вышеперечисленных институтов. Степень согласованности их ценностных посылов определяет уровень консолидации или раскола политической системы в целом\footnote{Соловьев А. И. Политология: Политическая теория, политические технологии: Учебник для студентов вузов. --- М.: Аспект Пресс, 2003. --- С. 355.}.